\subsection{Related Work}
\label{emt:sec:related}

\noindent
As memory translation has become a major bottleneck of emerging memory-hungry, irregular 
    workloads such as generative models, graph analytics,
    and recommendation systems,
    translation architecture has been an immensely active 
    research topic recently.
Extensive efforts are made on new designs of
    TLB~\cite{pham2012colt,pham2014increasing,park2017hybrid,barr2011spectlb} 
    and page walk caches~\cite{barr2010translation,bhattacharjee2013lmm},
    and more efficient huge page management~\cite{panwar2018making,navarro2002practical,kwon2016coordinated,panwar2019hawkeye,jia:pact:23}.

To fundamentally resolve translation bottlenecks, 
    recent work~\cite{skarlatos2020ecpt,kanellopoulos2023utopia,gupta2021midgard,yaniv2016hash,gosakan2023mosaic} 
    is actively rethinking translation architectures;
    a common thread is to prioritize speed over space efficiency (with abundant memory)
    and use fast lookup structures to reduce tail latency (avoiding pointer chasing).

% To reduce the memory translation overhead of emerging data-intensive applications,
% many new memory translation architectures have been proposed, including
% tree-based memory translation designs~\cite{park2022fpt,ainsworth2021compendia,
% ahn2012revisiting,margaritov2019asap,gandhi2016agile,alverti2020enhancing,self:dmt},
% hashing-based memory translation designs~\cite{jovan2022nestedecpt,skarlatos2020ecpt,
% gosakan2023mosaic},
% and many others~\cite{gupta2021midgard,kanellopoulos2023utopia,
% picorel2017near}.
%
Memory translation architectures have profound implications on OS reliability and performance,
    but unfortunately have not been well explored in prior work.
EMT is designed to enable development and evaluation of OS support for emerging translation architectures
    to formulate an OS perspective.

Note that memory management and its extensibility were well
    studied in OS research~\cite{krueger1993tools,rashid1987machine,abrossimov1989generic,engler1995avm,bershad1995extensibility}
    and are continuously revisited, driven by the growth 
    of memory capacity and the heterogeneity introduced by memory tiering~\cite{Tabatabai2024,Raybuck2021,Jalalian2024,li2023pond,maruf2023tpp,Bergman2022}.
So far, extensibility refers to the ability for user space 
    to customize memory management features, e.g., page fault handler.
Few prior studies considered extensibility in terms of empowering 
    memory translation architectures
    or studied their implications on OS performance.

%
% Most works mainly center on the design
% and implementation of the extensible memory management without focusing on extensible memory translation.
%

We are inspired by machine-independent memory management designs in 
    Mach~\cite{rashid1987machine} and SunOS~\cite{gingell1987sunos}.
%    and developed the EMT's interface on those ideas for Linux.
The goal of EMT
    is to enable the OS development for emerging, experimental
    memory translation architectures and 
    demonstrate the value with our experience.
We attempted to simplify the EMT interface, e.g., to a map interface~\cite{rashid1987machine,gingell1987sunos};
	however, we find it hard to balance architecture neutrality 
	and low-level optimizations.
% Unlike a map interface, the exposure of low-level detail in EMT allows such task to be implemented
%	in arch-neutral code with near-Linux performance.
% It also enables researchers to develop prototypes using existing Linux infrastructures,
%	without worrying about adding polyfills for uninterested architectures,
%	hence providing better support in the sense of functionalities.
One way is to add hints, but 
	% (e.g.,
 	% the last page found) and passes it down to the MMU drivers. \alan{This makes no sense to me}
	 hints are not expressive for passing information to different MMU drivers.

% \schai{Virtuoso is a very recent hardware simulation framework that 
% 	considers OS overhead... using a novel idea called mimic OS...
% 	what's different? 
% EMT is complementary as [[the goal is...]]
% EMT can do [[what it cannot do]]}

\schai{
% Virtual memory subsystem resarch involves innovation in both the OS software and hardware architecture.
% The existing simulators are either running short in accuracy of modeling software compoments
%     or too slow and complicated to evaluate the software-hardware co-design
%     ~\.
% Prototype and evaulation frameworks for new memory translation architectures
% Given extensive efforts on innovating translation architectures, TLBs, and
%     huge page memory management,
    %  introduces a simulatin framework
    % to enable quick and accurate prototyping and evaluation of virtual memory subystem.
Virtuoso~\cite{Kanellopoulos2025virtuoso} is a recent hardware simulation framework that 
	considers OS overhead (mostly the overhead of page fault handling) 
    using an imitated MimicOS in the user space.
% It balances the accuracy of modeling software components and the speed of simulation.
% It enables quick and accurate prototyping and evaluation of virtual memory subystem.
EMT and Virtuoso are fundamentally different but are complementary.
Unlike Virtuoso which is designed for architecture-centric simulations,
% which aims to avoid implementing the complexity of
%    full kernel implementation without compromising the simulation accuracy.
EMT is an OS interface for enabling OS developement and evaluation
    with regard to new translation architectures.
% In comparison, EMT is a OS-centric framework which aims to reduce the complexity of implementing OS support
%    by clearly defining the abstraction for translation logic and connecting the OS with hardware.
EMT benefits architecture designs from an OS perspective
    by running real MMU drivers on Linux.}
%    and model the OS performance impact more accurately (\S\ref{emt:sec:eval}).}